\chapter{Rechenregeln für reelle Folgen}
\label{b21guide:ch:calculus}

Alle Folgen in diesem Kapitel sind reell. Die Voraussetzungen
stehen jeweils beim Satz.

\section{Punktweise Operationen}

\begin{definition}[Rechenoperationen auf Folgen]
Summe, Differenz, Produkt, skalares Vielfaches und Betrag werden
gliedweise gebildet:
\[
 \begin{aligned}
 (a+b)_n&=a_n+b_n,& (a-b)_n&=a_n-b_n,\\
 (ab)_n&=a_nb_n,& (ca)_n&=ca_n,& (|a|)_n&=|a_n|.
 \end{aligned}
\]
Der Quotient hat die Glieder $a_n/b_n$ an denjenigen Indizes,
an denen $b_n\ne0$ ist. Eine Quotientenfolge auf ganz $\mathbb N$
verlangt $b_n\ne0$ für alle $n$.
Ist der Nenner nur schließlich ungleich $0$, wird die Formel
auf einer entsprechenden Restfolge verwendet; endlich viele
fehlende Anfangswerte dürfen beliebig reell ergänzt werden.
\end{definition}

\section{Algebraische Grenzwertregeln}

\begin{theorem}[Lineare Kombinationen und Produkte]
\label{b21guide:thm:algebra}
Aus $a_n\to L$ und $b_n\to K$ folgt für $c,d\in\mathbb R$
\[
 ca_n+db_n\to cL+dK,\qquad a_nb_n\to LK.
\]
Die Regeln gelten durch wiederholte Anwendung auch für jede
feste endliche Anzahl von Summanden oder Faktoren.
\end{theorem}

\begin{theorem}[Kehrwerte und Quotienten]
\label{b21guide:thm:quotient}
Gelte $a_n\to L$, $b_n\to K$ und $K\ne0$.
Dann ist $b_n$ schließlich ungleich $0$ und auf dieser Restfolge gilt
\[
 \frac1{b_n}\to\frac1K,\qquad
 \frac{a_n}{b_n}\to\frac LK.
\]
Bei überall von $0$ verschiedenen Nennern gelten die Aussagen
für die ganzen Quotientenfolgen.
\end{theorem}

\begin{theorem}[Nullfolge mal beschränkte Folge]
\label{b21guide:thm:null-bounded}
Ist $a_n\to0$ und $(b_n)$ beschränkt, so gilt $a_nb_n\to0$.
Konvergenz von $(b_n)$ ist dafür nicht erforderlich.
\end{theorem}

\begin{theorem}[Betrag, Minimum und Maximum zweier Folgen]
Aus $a_n\to L$ und $b_n\to K$ folgt
\[
 \begin{aligned}
 |a_n|&\to |L|,\\
 \min\{a_n,b_n\}&\to\min\{L,K\},\\
 \max\{a_n,b_n\}&\to\max\{L,K\}.
 \end{aligned}
\]
Insbesondere gilt $|a_n|\to0$ genau dann, wenn $a_n\to0$.
\end{theorem}

\begin{theorem}[Feste ganzzahlige Potenzen]
Sei $a_n\to L$. Für jede feste ganze Zahl $m\geq1$ gilt
$a_n^m\to L^m$.
Ist $L\ne0$, so gilt außerdem auf einer Restfolge
$a_n^{-m}\to L^{-m}$.
\end{theorem}

\begin{theorem}[Feste Wurzeln]
\label{b21guide:thm:roots}
Sei $r\geq1$ eine feste ganze Zahl und $a_n\to L$.
Für gerade $r$ sei zusätzlich $a_n\geq0$ für alle betrachteten
Indizes; dann ist auch $L\geq0$. Es gilt
\[
 \sqrt[r]{a_n}\longrightarrow\sqrt[r]{L}.
\]
Für ungerade $r$ gilt die Aussage für beliebige reelle $a_n,L$.
Bei geradem $r$ bezeichnet die Wurzel stets die nichtnegative Wurzel.
\end{theorem}

\Needspace{19\baselineskip}
\section{Ein Kriterium über aufeinanderfolgende Glieder}

\begin{theorem}[Quotientenkriterium für Folgen]
\label{b21guide:thm:ratio}
Es sei $a_n\ne0$ schließlich. Existiere der endliche Grenzwert
\[
 \left|\frac{a_{n+1}}{a_n}\right|\longrightarrow q.
\]
Dann ist $q\geq0$ und es gelten:
\begin{enumerate}
 \item Ist $q<1$, so ist $a_n\to0$.
 \item Ist $q>1$, so überschreitet $|a_n|$ schließlich
 jede reelle Zahl, also $|a_n|\to+\infty$ im Sinn von
 Definition~\ref{b21guide:def:infinity}.
\end{enumerate}
Der Fall $q=1$ entscheidet weder Konvergenz noch Divergenz.
Ohne Vorzeichenannahme folgt im zweiten Fall keine bestimmte
Richtung für $a_n$ selbst.
\end{theorem}

Hier handelt es sich um ein Resultat über die Folgenglieder,
nicht um ein Konvergenzkriterium für eine Reihe.

\section{Endliche Differenzen und polynomiale Folgen}

\begin{definition}[Vorwärtsdifferenzen]
\label{b21guide:def:differences}
Für eine reelle Folge $a$ werden ihre Differenzenfolgen
rekursiv festgelegt durch
\[
 (\Delta^0a)_n=a_n,\qquad
 (\Delta^{r+1}a)_n=(\Delta^ra)_{n+1}-(\Delta^ra)_n
 \quad(r,n\in\mathbb N).
\]
Insbesondere ist $(\Delta a)_n=a_{n+1}-a_n$.
Für die folgenden endlichen Summen verwenden wir
$\binom nk=n!/(k!(n-k)!)$ bei $0\leq k\leq n$ und
$\binom nk=0$ bei $k>n$; dabei ist $0!=1$ und
$n!=1\cdot2\cdots n$ für $n\geq1$.
\end{definition}

\begin{theorem}[Newtondarstellung und Abbruch der Differenzen]
\label{b21guide:thm:newton}
Für jede reelle Folge und jeden Index $n\in\mathbb N$ gilt
die endliche Identität
\[
 a_n=\sum_{k=0}^{n}\binom nk(\Delta^ka)_0.
\]
Für jede feste Zahl $d\in\mathbb N$ sind folgende Aussagen
äquivalent:
\begin{enumerate}
 \item $(\Delta^{d+1}a)_n=0$ für alle $n\in\mathbb N$.
 \item Es gibt reelle Zahlen $c_0,\ldots,c_d$ mit
 \[
  a_n=c_0+\sum_{j=1}^{d}c_jn^j
  \qquad(n\in\mathbb N).
 \]
 \item Für alle $n\in\mathbb N$ gilt
 \[
  a_n=\sum_{k=0}^{d}\binom nk(\Delta^ka)_0.
 \]
\end{enumerate}
Eine Folge mit diesen Eigenschaften heißt \emph{polynomiale
Folge vom Grad höchstens $d$}. Die Koeffizienten
$c_0,\ldots,c_d$ sind eindeutig bestimmt.
Unter dieser Voraussetzung bestimmen die ersten $d+1$ Glieder
bereits die gesamte Folge.
Für $d=0$ bedeutet dies genau, daß die Folge konstant ist.

Das Verschwinden der Differenzen nur auf einem endlichen
Anfangsstück charakterisiert noch keine polynomiale Folge
auf ganz $\mathbb N$.
\end{theorem}

Separate Ergänzungen:
\href{../newton/_B21-newton-reading.pdf\#newton.reading}{Lesefassung}
und \href{../newton/_B21-newton-proofs.pdf\#newton.proofs}{Beweistabellen}.

\clearpage
\section{Endliche Summationsformeln}

\begin{theorem}[Teleskopformel und Abelsche partielle Summation]
\label{b21guide:thm:finite-summation}
Für eine reelle Folge $(u_n)$ und natürliche Zahlen $r\leq s$ gilt
die \emph{Teleskopformel}
\[
 \sum_{k=r}^{s}(u_{k+1}-u_k)=u_{s+1}-u_r.
\]

Für reelle Folgen $(a_n),(b_n)$ setze
$A_j=\sum_{k=r}^{j}a_k$ für $r\leq j\leq s$.
Dann gilt die \emph{Abelsche Summationsformel}
\[
 \sum_{k=r}^{s}a_kb_k
 =A_sb_s+\sum_{k=r}^{s-1}A_k(b_k-b_{k+1}).
\]
Ist zusätzlich $|A_j|\leq C$ für alle $r\leq j\leq s$
mit einem $C\geq0$ und gilt
$b_r\geq b_{r+1}\geq\cdots\geq b_s\geq0$, so folgt die
endliche Abschätzung
\[
 \left|\sum_{k=r}^{s}a_kb_k\right|\leq Cb_r.
\]
Für $r=s$ ist die Summe von $r$ bis $s-1$ leer und gleich $0$.
Alle Aussagen betreffen nur endliche Indexabschnitte;
Konvergenzannahmen werden nicht benötigt.
\end{theorem}

Separate Ergänzung: \href{../abel/_B21-abel-reading.pdf\#abel.reading}{Lesefassung mit kurzen Beweisen, Beispielen und historischen Originalquellen}.
